\section{Why This Chapter Must Become Stricter}

The earlier version of this chapter correctly saw that frameworks are formed,
not merely encountered. But it still read too much like a mechanism sketch. If
the manuscript wants to claim that habits, interpretive frames, and behavioural
regimes are historically produced structures, then it must explain
\emph{how} higher-level organization arises, \emph{what} kind of evidence
supports that claim, and \emph{where} the analogy to human behaviour becomes
interpretive rather than directly measured.

This chapter therefore adopts a visible three-layer architecture:
\begin{enumerate}[nosep]
  \item \textbf{Biological grounding:} what current biology and neuroscience
        allow us to say about organized state formation and stability.
  \item \textbf{Physical and systems analogies:} what complex-systems and
        self-organization language contributes when used carefully.
  \item \textbf{Mathematical and explanatory structure:} what emergence,
        path dependence, metastability, and multiscale organization mean in a
        disciplined model.
\end{enumerate}

The evidence policy remains explicit. The chapter's main direct sources are
\emph{Metastability demystified: the foundational past, the pragmatic present
and the promising future} (Nature Reviews Neuroscience, 2024, \textbf{medium}),
\emph{Associative conditioning in gene regulatory network models increases
integrative causal emergence} (Communications Biology, 2025, \textbf{medium}),
and the neighbouring empirical consciousness chapter, whose strong sources
justify taking state structure, stability, and transition seriously at the
biological level. The chapter also leans on classic textbooks such as John
Holland's \emph{Emergence}, \emph{Introduction to the Theory of Complex
Systems}, and mathematically oriented dynamical-systems texts. None of these
sources directly proves the full behavioural framework of this book. Their
value is to constrain the explanation so it is structurally serious rather than
metaphorical improvisation.
